\documentclass{article}

\title{Relazione 3}

\begin{document}
	
	\maketitle
	
	
	
	\section{Strumenti di misura e approccio sperimentale}
	
	Nella presente esperienza vengono misurati gli ingrandimenti visuali di un microscopio semplice e di un microscopio composto.
	Gli angoli visuali vengono misurati utilizzando un modello d'occhio realizzato attraverso una lente convergente che funge da cristallino, con una focale nominale di 150mm, ed uno schermo traslucido che funge da retina.
	Il microscopio semplice è costituito da una lente convergente con un focale nominale di 100mm, da utilizzare come lente di ingrandimento. Tale lente sarà successivamente utilizzata come oculare del microscopio composto e le verrà anteposta una lente convergente con una focale nominale di 50mm come obiettivo.
	Le lenti e lo schermo sono montate su un banco ottico tramite dei cavalieri, i quali vi possono scorrere per mezzo di un binario. Sul banco ottico, di lunghezza circa un metro, è inoltre presente una scala millimetrata, che permette di misurare le distanze relative degli elementi del sistema ottico. In particolare, essendo i cavalieri spessi circa 5cm, la distanza tra due elementi verrà misurata come la differenza delle posizioni dei bordi corrispondenti dei rispettivi cavalieri alla quale sarà associata un'incertezza massima di 5mm, pari allo spessore dei suddetti cavalieri.
	Per misurare gli angoli visuali si utilizza un oggetto di proiezione di cui si misura la dimensione di un dettaglio, da confrontare di volta in volta con le dimensione del dettaglio nell'immagine virtuale prodotta dalle lenti. 
	L'oggetto viene montato sul primo cavaliere, partendo dall'origine del banco ottico.
	Prima di eseguire le misure, per ogni configurazione  ci si assicura che i centri dell'oggetto e delle lenti siano allineati e che i piani in cui essi giacciono siano paralleli e ortogonali alla direzione del banco ottico. 
	Il primo ingrandimento visuale da misurare è quello del microscopio semplice, per cui, innanzitutto, si sceglie un particolare dell'oggetto di proiezione e si misura la sua dimensione con un regolo di sensibilità $1mm. 
	Preliminarmente si misura l'angolo visuale sotto il quale si vede il dettaglio dell'oggetto quando l'occhio è accomodato a una visione riposante, ovvero quando la distanza tra il cristallino e l'oggetto è la distanza di visione distinta $d_0=(25,0\pm0,5)cm$. Quindi si fissa la distanza oggetto-cristallino a $d_0$ e si monta la retina dalla parte opposta del cristallino rispetto all'oggetto, spostando quest'ultima fino a raccogliere l'immagine a fuoco. Si osserva che esistono due posizioni $q_r^{min}$ e $q_r^{max}$ tra le quali l'immagine risulta a fuoco, in corrispondenza delle quali la dimensione del dettaglio dell'immagine è rispettivamente $Y_2^{min}$ e $Y_2^{max}$. Pertanto, si misurerà $Y_2$ come media di $Y_2^{min}$ e $Y_2^{max}$ più o meno la loro semidifferenza e parimenti si farà per $q_r$.
	L'angolo $\alpha_r$ è piccolo e dunque circa uguale alla sua tangente, ovvero al rapporto tra $Y_2$ e $q_r$.
Successivamente si monta la lente d'ingrandimento a una distanza $p$ minore della focale dall'oggetto e si pone il cristallino nel secondo fuoco della lente. In questa configurazione l'ingrandimento visuale è indipendente da $p$ e vale $G_v=\frac{d_0}{f}$. Quindi si misurano $q_2$ e $Y_2$ con il procedimento già utilizzato nella misura di $alpha_r$. In  questo modo si può stimare l'angolo visuale attraverso la sua tangente, in quanto è piccolo: $\beta=\frac{Y_2}{q_2}$. Le misure vengono ripetute per diversi valori di $p$, in modo da ottenere diverse stime dell'ingrandimento visuale $G_v=\frac{\beta}{\alpha_r}$ (vedere Tabella 1). Quindi si riporta su un grafico $G_v$ in funzione di $p$: come atteso le stime di $G_v$ sono compatibili entro gli errori sperimentali. Si può eseguire un fit lineare per cercare un'eventuale dipendenza dell'ingrandimento lineare da $p$: la stima del coefficiente angolare è compatibile con 0, dunque si può combinare le diverse stime con una media pesata ottenendo la stima finale $G_v=$.
Per la misura dell'ingrandimento visuale del microscopio composto, si inizia scegliendo un altro particolare dell'oggetto, di cui si misura la dimensione, e per il quale si ripete la stima dell'angolo $alpha_r^'$, attraverso la stima di $Z_r$ e $Q_r$. Mantenendo fissa la distanza $Q_r$ tra cristallino e retina, si montano l'obiettivo e l'oculare in modo che il cristallino sia nel secondo fuoco dell'oculare e che il tiraggio meccanico sia $T=(35,0\pm0,5)cm$. Quindi si varia la distanza $p$ dell'oggetto dall'oculare per valori maggiori dellla focale, fino a ottenere un'immagine nitida sulla retina. Misurando $Z_2$, si stima $\beta^'$ e con questo $G_v^'$. Come atteso la stima è compatibile con il valore teorico $G_v^'=\frac{\Delta d_0}{f_{ob} f_{oc}}.
	

	
	
		\section{conclusioni}

			
			
			

			
			
			
\end{document}